\chapter{Earache}

\section*{Definition}
Pain or discomfort localized to the ear, ranging from dull ache to sharp stabbing sensation, originating from the external, middle, or inner ear structures or referred from adjacent regions.

\section*{What Happens in Your Body}
Ear pain arises from inflammation, infection, pressure changes, or irritation of the auricle, ear canal, tympanic membrane, or middle ear structures. Referred pain occurs via shared innervation with the temporomandibular joint, teeth, sinuses, and pharynx through cranial nerves V, VII, IX, and X.

\section*{Causes}
\begin{itemize}
  \item Otitis media (middle ear infection)
  \item Otitis externa (swimmer's ear)
  \item Eustachian tube dysfunction
  \item Temporomandibular joint (TMJ) disorders
  \item Dental infections or abscesses
  \item Foreign body in ear canal
  \item Earwax impaction
  \item Perforated tympanic membrane
  \item Mastoiditis
  \item Referred pain from tonsillitis or pharyngitis
\end{itemize}

\section*{When to See a Doctor}
See your doctor if:
\begin{itemize}
  \item Symptoms are severe or getting worse over time
  \item Severe ear pain, fever, hearing loss, ear discharge, or pain persisting more than 2-3 days
  \item You have other concerning symptoms such as fever, weight loss, or bleeding
\end{itemize}

\section*{Related Symptoms}
\begin{itemize}
  \item Hearing loss
  \item Fever
  \item Ear discharge
  \item Dizziness
  \item Headache
\end{itemize}
\textbf{Important:} If this symptom persists, your doctor will check for serious underlying causes.

\section*{Self-Care / Home Management}
\begin{itemize}
  \item Rest and avoid activities that make it worse.
  \item Maintain adequate hydration and a balanced diet.
  \item Over-the-counter remedies may provide symptomatic relief (consult a pharmacist or doctor).
  \item Monitor symptoms and seek professional advice if they persist or worsen.
  \item Avoid self-medicating with prescription medications without medical guidance.
\end{itemize}

\section*{How Your Doctor Will Evaluate You}
\begin{itemize}
  \item A thorough history and physical examination are the first steps in evaluation.
  \item Laboratory tests (blood work, urinalysis) may be ordered based on clinical suspicion.
  \item Imaging studies (X-ray, ultrasound, CT, MRI) may be indicated when structural pathology is suspected.
  \item Specialist referral is arranged when the diagnosis remains uncertain or requires specific expertise.
\end{itemize}

\section*{Treatment Options}
\begin{itemize}
  \item Treatment targets the underlying cause identified through diagnostic evaluation.
  \item Supportive care and symptom management are provided while awaiting definitive therapy.
  \item Medications, lifestyle modifications, and physical therapies may be prescribed as appropriate.
  \item Follow-up ensures treatment effectiveness and allows timely adjustments to the care plan.
\end{itemize}

\clearpage